\chapter{The Density Question, Resolved Sharply}
\label{ch:density}

Chapter~\ref{ch:minimal-witnesses}'s fourth open question asked whether,
for a random context graph, the probability that admissible edge data
requires $\C$ approaches $1$ as edge density increases — an
Erd\H{o}s--R\'enyi-style asymptotic threshold, explicitly flagged there as
needing a probability measure on edge-label data this book had not yet
defined. This chapter defines that measure and answers the question, but
the answer is sharper than the question anticipated: there is no
asymptotic threshold to find, because the threshold is already crossed at
the first cycle.

\section{A Natural Measure, and a Correction}

\begin{definition}[Independent continuous edge measure]
\label{def:edge-measure}
\index{Independent continuous edge measure (definition)}
For a context graph $G$ at $d=2$, equip each edge with an independent
$\mathrm{Uniform}[0,\pi/4]$ distribution over its canonical angle,
independently across edges.
\end{definition}

This is the natural candidate the question calls for: no correlation
between edges is assumed, and no edge value is privileged.

Before proceeding, a claim made in passing in
Chapter~\ref{ch:real-mub-obstruction} needs correcting. That chapter's
discussion of the triangle-gluing theorem stated that "generic triples of
edge angles do admit real joint realizations." Checked directly against
Definition~\ref{def:edge-measure}, this is false, and has been corrected
at the point it was made (Chapter~\ref{ch:real-mub-obstruction}'s remark
following Theorem~\ref{thm:triangle-gluing}): realizability is not the
generic case. It is the exception.

\section{The Sharp Result}

\begin{theorem}[Realizability is measure-zero once $\beta_1(G)\ge1$]
\label{thm:density-sharp}
\index{Realizability is measure-zero once beta1(G)>=1 (theorem)}
For any context graph $G$ with $\beta_1(G)\ge1$, under
Definition~\ref{def:edge-measure}, $\Pr[\text{edge data is globally
real-realizable}] = 0$.
\end{theorem}

\begin{proof}
Fix a spanning tree $T$ and the $\beta_1(G)$ non-tree edges. By Theorem
\ref{thm:graph-realizability}, global realizability requires a single
choice of tree-edge lifts under which \emph{every} non-tree edge matches
its target angle. For fixed tree-edge angles, each non-tree edge's target
must equal one of finitely many values determined by the tree data
(Theorem~\ref{thm:triangle-gluing} for each fundamental cycle) — a
codimension-one condition on that edge's independently-drawn value. With
$\beta_1(G)$ non-tree edges drawn independently, satisfying all
$\beta_1(G)\ge1$ such conditions simultaneously is a codimension-$\beta_1(G)$
event under a continuous product measure, hence has probability zero.
\end{proof}

This was verified computationally before being stated as a theorem, at
two different values of $\beta_1$: $0.05\%$ of $2{,}000{,}000$ fully
independent triangle trials matched realizability within a generous
numerical tolerance of $10^{-4}$ (consistent with probability exactly
zero as the tolerance shrinks), and $0$ of $200{,}000$ trials matched on
$K_4$ ($\beta_1=3$), the higher codimension making the event
correspondingly harder to hit even at finite tolerance.

\section{Why This Is Not an Erd\H{o}s--R\'enyi Threshold}

The question as posed in Chapter~\ref{ch:minimal-witnesses} anticipated an
asymptotic phenomenon: a critical edge density below which real
realizability is typical and above which it is not, in the manner of the
giant-component threshold. Theorem~\ref{thm:density-sharp} shows this
picture does not apply. There is no gradual transition and no critical
density to locate: the moment a graph contains even one cycle
($\beta_1(G)=1$, the minimal case, achieved by a bare triangle), generic
edge data already fails to realize over $\R$ with probability one. Adding
further edges (increasing $\beta_1(G)$ further) does not create a
threshold that was previously absent; it was already crossed.

\begin{remark}
This should be read alongside Chapter~\ref{ch:higher-d-witnesses}'s
inversion result in the same spirit: an intuition imported from a
different context (the Erd\H{o}s--R\'enyi threshold, or "higher
dimension means larger witnesses") turns out not to describe what
actually happens here, and the correction is sharper and more interesting
than the naive expectation. Where Chapter~\ref{ch:higher-d-witnesses}
found the minimal witness \emph{shrinks} rather than grows with dimension,
this chapter finds the density threshold \emph{degenerates to a single
point} ($\beta_1=1$) rather than sweeping out a nontrivial range.
\end{remark}

\begin{ledgerentry}[End of Chapter~\ref{ch:density}]
\textbf{Established:} the density-threshold question of Chapter
\ref{ch:minimal-witnesses} is resolved (Theorem~\ref{thm:density-sharp}):
under the natural independent-continuous edge measure, real realizability
has probability zero for every graph with $\beta_1(G)\ge1$, not an
asymptotic phenomenon requiring a critical density. A false claim in
Chapter~\ref{ch:real-mub-obstruction} ("generic triples do admit real
realizations") is corrected at its source as part of this resolution.\\
\textbf{Not established:} anything about non-continuous or correlated
measures on edge data, which might in principle behave differently (e.g.\
a measure concentrated on real-derived triples would trivially give
probability one, the opposite regime); anything about $d\ge3$, where
Theorem~\ref{thm:density-sharp}'s proof — via the fundamental-cycle
characterization of Theorem~\ref{thm:graph-realizability} — still applies
verbatim to the graph-level obstruction, but where Chapter
\ref{ch:higher-d-witnesses}'s cheaper single-edge obstruction is a
separate phenomenon this chapter's measure has not been extended to
analyze.
\end{ledgerentry}
